\section{Conclusion}
\label{sec:conclusion}

We presented SCTW and Adjoint-SCTW, two training-free population meshes for Euler sampling of
frozen flow-matching models. SCTW estimates the local material acceleration along a fine pilot.
Adjoint-SCTW further weights that defect by terminal-state sensitivity using a constant-memory
randomized discrete adjoint. Both then place knots by a declared tempered equidistribution rule.

The strongest evidence is staged-shapes, where Adjoint-SCTW improves both learned-ODE integration
and endpoint feature W1 over SCTW, uniform Euler, and every tested fixed Euler mesh at each NFE.
The same terminal weighting substantially reduces learned-ODE integration error on GM16 and hard
Burgers. The CIFAR result defines the boundary: SCTW improves FID at moderate NFE, but
Adjoint-SCTW's more accurate integration does not consistently improve FID. Across schedules,
integration error and endpoint quality are nearly rank-aligned on staged-shapes and weakly or
negatively aligned on CIFAR at moderate and high NFE. A sampler cannot remove endpoint bias in the
learned vector field, and discretization error can sometimes cancel part of that bias.

The methods are therefore best understood as interpretable mesh diagnostics, not universally
optimal solvers. The remaining limitations are concrete: one population mesh may hide
sample-specific stiffness; separately norming transported defects ignores cancellation across
steps; randomized calibration adds an offline cost; and the tempering exponent remains declared
rather than derived. Multi-checkpoint image studies and comparisons with learned solver schedules
are needed before making broader generative-quality claims.
